\section{Conclusion}

In this article, a novel zero-shot fault diagnosis framework named the CLIP-Guided Dual-Path Diffusion Model (CDDM) provides a robust solution to the persistent challenge of unseen fault detection in complex industrial environments. Traditional generative zero-shot models heavily rely on manually annotated, discrete binary attributes, which invariably result in the loss of fine-grained semantic information and often suffer from severe issues such as mode collapse and domain shift during adversarial training. To overcome these critical bottlenecks, the proposed CDDM pioneers a cross-modal paradigm that explicitly bridges textual diagnostic descriptions and complex industrial time-series signals. By leveraging large language models (LLMs) alongside contrastive language-image pre-training (CLIP) architectures, textual fault causes and functional effects are mapped into a dense, continuous semantic space, completely eliminating the conventional reliance on hard-coded attribute matrices.

A major breakthrough of the CDDM framework lies in its generative dual-path diffusion architecture. Comprising a spatial-temporal path utilizing one-dimensional convolutional neural networks and a time-frequency spectral path bolstered by advanced attention mechanisms, the proposed model ensures stable, robust, and highly faithful signal generation. Furthermore, the integration of a multi-modal contrastive consistency loss guarantees that the generated unseen signal representations precisely align with the guiding text semantics. This structural innovation directly replaces standard attribute regression techniques and substantially mitigates the persistent domain shift problem. Extensive experimental benchmarks across diverse industrial datasets—including the Three-Phase Transmission System (TPTS), the Tennessee Eastman Process (TEP), and a complex hydraulic system—demonstrate that CDDM achieves state-of-the-art zero-shot diagnostic performance, consistently outperforming existing generative adversarial network (GAN) and variational autoencoder (VAE) baselines in both sample distribution fidelity and classification accuracy \cite{tzeng2017adversarial}.

The industrial implications of this research are highly significant for modern Industry 4.0 applications \cite{wang2020deep,lee2014prognostics}. In safety-critical sectors where certain catastrophic faults are inherently too dangerous, rare, or costly to reproduce, the ability to synthesize physically plausible signals strictly from expert textual descriptions enables the training of highly comprehensive predictive maintenance systems \cite{lee2014prognostics}. This approach drastically minimizes the dependency on extensive historical fault datasets, thereby enhancing the functional resilience, scalability, and operational safety of modern machinery and continuous chemical processes \cite{precup2015overview}.

Despite these substantial advancements, there remain certain constraints to be addressed in future research. The inherent iterative nature of the generative diffusion process introduces computational latency during the reverse denoising phase, which may hinder its immediate deployment in strictly real-time diagnostic paradigms on resource-constrained edge devices. Future scopes will therefore focus on implementing accelerated sampling strategies, such as denoising diffusion implicit models (DDIM \cite{song2021denoising}) and knowledge distillation, to optimize sample generation efficiency. Additionally, extending the proposed dual-path conditioning framework to accommodate and fuse broader multi-modal sensor arrays—such as transient acoustic emissions and full-field thermal imaging—promises to further elevate the robustness and real-world applicability of generalized zero-shot condition monitoring systems.
